\documentclass[10pt]{beamer}

\usepackage[spanish]{babel}
\usepackage[utf8]{inputenc}
\usetheme{metropolis}

\usepackage{appendixnumberbeamer}
\usepackage{pgfplots}
\usepackage{booktabs}
\usepackage[scale=2]{ccicons}
\usepackage{minted}
\usepgfplotslibrary{dateplot}

\usepackage{xspace}
\newcommand{\themename}{\textbf{\textsc{metropolis}}\xspace}

\usemintedstyle{emacs}

\newcommand\Fontvi{\fontsize{10}{12}\selectfont}


\title{Property-Based Testing}
% \subtitle{A modern beamer theme}
\date{\today}
\author{
  Luis Eduardo Bueso de Barrio\\
  lbueso@acm.org
}
\graphicspath{{./data/}{../resources/}}
\institute{ACM PLIG}
\titlegraphic{
  \hfill\includegraphics[height=1.5cm]{pics/logo_acm_cromo.png}
  \includegraphics[height=1.5cm]{pics/sig_plig.png}
}

\begin{document}

\nocite{*}

\maketitle

\begin{frame}{Índice}
  \setbeamertemplate{section in toc}[sections numbered]
  \tableofcontents[hideallsubsections]
\end{frame}

\section{Introducción}

\begin{frame}                   %TODO
  El testing es el método de verificación de la calidad de software
  más usado.

  Problema: consume gran cantidad de recursos (hasta el 50\% del coste
  del software).

  Preguntas previas:
  \begin{itemize}
  \item ¿Quién prueba su código?
  \item ¿Quién realiza \textit{suites} de tests?
  \item ¿Quien realiza tests con herramientas de \textit{Property-Based Testing}?
  \item ¿Quién verifica formalmente su software?
  \end{itemize}
\end{frame}

\section{¿Qué es el Property-Based Testing?}

\begin{frame}{Definición}
  El \textit{Property-Based Testing} es una metodología de testing que
  a través de la especificación de propiedades, el análisis del flujo
  de datos del programa y la ejecución de tests busca su corrección y
  completitud.

  Responde a las siguientes preguntas:
  \begin{itemize}
  \item ¿Hasta dónde se puede llegar mediante el testing?
  \item ¿Qué datos se deben usar en los tests?
  \item ¿Cuándo se ha hecho suficiente testing?
  \item ¿Cómo se determina si un test es exitoso o fallido?
  \end{itemize}

\end{frame}

\begin{frame}{¿Qué ha conseguido hasta ahora?}
  El \textit{Property-Based Testing} ha evolucionado mucho desde su
  concepción, y ahora es una metodología muy madura, que reduce en
  gran medida el precio de la generación de tests.

  Ahora mismo existen gran cantidad de herramientas que utilizan esta
  metodología:
  \begin{itemize}
  \item Haskell \textit{QuickCheck}.
  \item Erlang \textit{QuickCheck}.
  \item Erlang \textit{PropEr}.
  \item Scala \textit{ScalaCheck}.
  \item Python \textit{Hypothesis}.
  \item Java 8 \textit{QuickTheories}.
  \item JavaScript \textit{JSVerify}.
  \end{itemize}
\end{frame}

\begin{frame}{¿Qué ha conseguido hasta ahora?}
  El \textit{Property-Based Testing} ha evolucionado mucho desde su
  concepción, y ahora es una metodología muy madura, que reduce en
  gran medida el precio de la generación de tests.

  Ahora mismo existen gran cantidad de herramientas que utilizan esta
  metodología:
  % \begin{columns}
  %   \begin{column}{0.5\textwidth}
  \begin{itemize}
  \item \textbf{Haskell \textit{QuickCheck}}.
  \item Erlang \textit{QuickCheck}.
  \item Erlang \textit{PropEr}.
  \item Scala \textit{ScalaCheck}.
  \item Python \textit{Hypothesis}.
  \item Java 8 \textit{QuickTheories}.
  \item JavaScript \textit{JSVerify}.
  \end{itemize}
  % \end{column}
  % \begin{column}{0.5\textwidth}
  %     %   \includegraphics[height=1.5cm]{pics/haskell_logo_star.jpg}
  % \end{column}
  % \end{columns}
\end{frame}

\section{QuickCheck}

\begin{frame}{¿Qué es QuickCheck?}
  \textit{QuickCheck} es una herramienta de creación y testeo de
  propiedades sobre los programas (\textit{Property-Based Testing}).

  Permite la descripción de propiedades mediante funciones en Haskell
  que se prueban sobre una entrada aleatoria.

  Permite la creación de generadores propios.

  \textit{QuickCheck} es un \textit{DSL} (\textit{Domain-Specific
    Language}) de especificación de tests.

  Herramienta muy ligera, la implementación original tenía sobre 300
  líneas de código.
\end{frame}

\begin{frame}{Propiedades}
  Ejemplo:
  \inputminted{text}{data/reverse.hs}
  Propiedades:
  \inputminted{haskell}{data/reverse_prop.hs}
\end{frame}


\begin{frame}[standout]
  \Fontvi
  \inputminted{text}{data/pruebas_rev.txt}
\end{frame}

% \begin{frame}[standout]
%   \textexclamdown Atención!
% \end{standout}

\begin{frame}[standout]
  \Fontvi
  \inputminted{text}{data/reverse_prop_pruebas_verb.txt}
\end{frame}

\begin{frame}[standout]
  \Fontvi
  \inputminted{text}{data/reverse_prop_pruebas_verb1.txt}
\end{frame}

\begin{frame}{Propiedades de funciones}
  Propiedad commutativa y asociativa del operador de composición
  \inputminted{haskell}{data/func_prop.hs}
\end{frame}

\begin{frame}[standout]
  \Fontvi
  \inputminted{text}{data/pruebas_func.txt}
\end{frame}

\begin{frame}{Condicionales}
  Operación de implicación $(==>)$ para condicionales

  Ejemplo:
  \inputminted{haskell}{data/max_prop.hs}
\end{frame}

\begin{frame}[standout]
  \Fontvi
  \inputminted{text}{data/pruebas_max_prop.txt}
\end{frame}

\begin{frame}{Monitorización}
  También podemos monitorizar datos de los tests:
  \inputminted{haskell}{data/monitor_test.hs}
\end{frame}

\begin{frame}[standout]
  \Fontvi
  \inputminted{text}{data/monitor_test_pruebas.txt}
\end{frame}

\begin{frame}{Estructuras Infinitas}
  Hay una sutil diferencia entre estas dos propiedades:
  \inputminted{haskell}{data/infinite.hs}
\end{frame}

\begin{frame}[standout]
  \Fontvi
  \inputminted{text}{data/infinite_prueba.txt}
\end{frame}

\begin{frame}{Generadores}
  Además de los generadores propios de QuickCheck se pueden
  implementar generadores propios.

  Para implementar generadores se
  utiliza la \textit{typeclass}:
  \inputminted{haskell}{data/arbitrary.hs}

  Ejemplo:
  \inputminted{haskell}{data/colours.hs}
\end{frame}

\begin{frame}[standout]
  \Fontvi
  \inputminted{text}{data/colours_prueba.txt}
\end{frame}

\begin{frame}{Generadores}
  Podemos mejorar nuestros generadores controlando, por ejemplo, la
  frecuencia con la que se generan determinados valores.

  \inputminted{haskell}{data/colours1.hs}
\end{frame}

\begin{frame}[standout]
  \Fontvi
  \inputminted{text}{data/colours1_prueba.txt}
\end{frame}

\begin{frame}{Modificadores}
  Podemos mejorar nuestros generadores aplicándoles modificadores, que
  aportan ciertas propiedades:

  \inputminted{haskell}{data/modificadores.hs}

  Para buscar más modificadores consultar $Test.QuickCheck.Modifiers$
\end{frame}

\begin{frame}[standout]
  \Fontvi
  \inputminted{text}{data/modificadores_prueba.txt}
\end{frame}

\begin{frame}{Modificadores}
  Con modificadores podemos conseguir que siempre se generen datos
  correctos para propiedades que requieren condiciones.

  Ejemplo:
  \inputminted{haskell}{data/modificadores2.hs}
\end{frame}

\begin{frame}[standout]
  \Fontvi
  \inputminted{text}{data/modificadores2_pruebas.txt}
\end{frame}

\section{Conclusiones y observaciones}

\begin{frame}
  Problemas:
  \begin{itemize}
  \item Definir propiedades en muchos casos, no es una tarea trivial.
  \item Los \textit{side-effects} pueden ser difíciles de verificar.
  \end{itemize}
  A pesar de estos problemas, definir propiedades, nos aporta una gran
  recompensa.

  Los lenguajes funcionales puros son más fáciles de testear.

  Algunos ejemplos reales de uso de \textit{QuickCheck}:
  \begin{itemize}
  \item Unificación.
  \item Propiedades en circuitos con \textit{Lava}.
  \item \textit{Pretty Printing}.
  \item Edison.
  \end{itemize}
\end{frame}

\appendix

\begin{frame}{Referencias}

  \bibliographystyle{acm}
  \bibliography{presentacion}

\end{frame}

\begin{frame}[standout]
  ¿Preguntas?
\end{frame}

\end{document}